\documentclass[11pt,a4paper,oneside]{scrartcl}
\usepackage[utf8]{inputenc}
\usepackage[ngerman]{babel}
\usepackage{microtype}
\usepackage{booktabs}
\usepackage{enumitem}
\usepackage{tcolorbox}
\usepackage{siunitx}

\usepackage[pdftex, hidelinks]{hyperref}
\sisetup{range-phrase = {--}, range-units = single}
\newcommand{\celsius}[1]{\qty{#1}{\celsius}}
\newcommand{\timelineField}[2]{%
    \Form{\TextField[
        name=#1, width=7.5em, align=1, readonly=true,
        calculate={var eT=this.getField("endTime").value; var eD=this.getField("endDate").value; if(eT&&eD&&eT.indexOf(":")!==-1){var p=eT.split(":"); var eH=parseInt(p.shift(),10); var eM=parseInt(p.shift(),10); var dParts=eD.split("."); var dD=parseInt(dParts.shift(),10); var dM=parseInt(dParts.shift(),10)-1; var dY=parseInt(dParts.shift(),10); if(!isNaN(eH)&&!isNaN(eM)&&!isNaN(dD)&&!isNaN(dM)&&!isNaN(dY)){var date=new Date(dY,dM,dD,eH,eM,0,0); date.setMinutes(date.getMinutes()-#2); var curD=date.getDate(); var curM=date.getMonth()+1; var cH=date.getHours(); var cM=date.getMinutes(); var strDate=((curD<10)?"0"+curD:curD)+"."+((curM<10)?"0"+curM:curM)+"."; var strTime=((cH<10)?"0"+cH:cH)+":"+((cM<10)?"0"+cM:cM); event.value=strDate+" "+strTime;}}else{event.value="--.--. --:--";}}
    ]{}}%
}

\usepackage{auxhook}
% ingredients.tex
\usepackage{xfp} % Modern high-precision floating point calculation engine
\usepackage{tcolorbox}
\usepackage{booktabs}

% --- LATEX3 KEY-VALUE SETUP & MATH ENGINE ---
\ExplSyntaxOn

% Define the nutrition key storage fields
\keys_define:nn { ing }
  {
    name    .tl_set:N  = \l_ing_name_tl,
    kcal    .fp_set:N  = \l_ing_kcal_fp,
    protein .fp_set:N  = \l_ing_protein_fp,
    carbs   .fp_set:N  = \l_ing_carbs_fp,
    fat     .fp_set:N  = \l_ing_fat_fp,
    satfat  .fp_set:N  = \l_ing_satfat_fp,
    sugar   .fp_set:N  = \l_ing_sugar_fp,
    fiber   .fp_set:N  = \l_ing_fiber_fp,
    salt    .fp_set:N  = \l_ing_salt_fp,
  }

% Native LaTeX3 definition of the print macro to prevent parameter bugs
\cs_new:Npn \printValue #1 
  {
    \fpeval{round(#1, 2)}
  }

% Native LaTeX3 initialization of calculation registers
\cs_new:Npn \initRecipeCalculator 
  {
    \gdef\totalKcal{0}
    \gdef\totalProtein{0}
    \gdef\totalCarbs{0}
    \gdef\totalFat{0}
    \gdef\totalSatFat{0}
    \gdef\totalSugar{0}
    \gdef\totalFiber{0}
    \gdef\totalSalt{0}
  }

% Native LaTeX3 definition to save ingredients into the global data cache
\cs_new:Npn \defIngredient #1 #2
  {
    \expandafter\gdef\csname ingdata@#1\endcsname{#2}
  }

% Core logic to add an ingredient, evaluate math and print layout columns
\newcommand{\addIngredient}[2]{
    % Reset variables to default 0 values safely
    \keys_set:nn { ing } { name={#1}, kcal=0, protein=0, carbs=0, fat=0, satfat=0, sugar=0, fiber=0, salt=0 }
    % Load and parse the stored key-value list for this specific ID
    \keys_set:nv { ing } { ingdata@#1 }
    %
    % Export the token list name safely to a standard global macro so tabular can read it
    \tl_gset_eq:NN \g_printed_name_tl \l_ing_name_tl
    \expandafter\gdef\expandafter\printedName\expandafter{\g_printed_name_tl}
    %
    % Accumulate values using floating point evaluation
    \xdef\totalKcal{\fpeval{\totalKcal + (#2 * \l_ing_kcal_fp / 100)}}
    \xdef\totalProtein{\fpeval{\totalProtein + (#2 * \l_ing_protein_fp / 100)}}
    \xdef\totalCarbs{\fpeval{\totalCarbs + (#2 * \l_ing_carbs_fp / 100)}}
    \xdef\totalFat{\fpeval{\totalFat + (#2 * \l_ing_fat_fp / 100)}}
    \xdef\totalSatFat{\fpeval{\totalSatFat + (#2 * \l_ing_satfat_fp / 100)}}
    \xdef\totalSugar{\fpeval{\totalSugar + (#2 * \l_ing_sugar_fp / 100)}}
    \xdef\totalFiber{\fpeval{\totalFiber + (#2 * \l_ing_fiber_fp / 100)}}
    \xdef\totalSalt{\fpeval{\totalSalt + (#2 * \l_ing_salt_fp / 100)}}
    %
    % Output table row with exactly 2 decimal places
    \fpeval{round(#2, 2)} & g & \printedName \\
}
\ExplSyntaxOff

% --- NUTRITION FACTS PANEL GRAPHIC ---
\newcommand{\printNutritionFacts}{%
    \begin{tcolorbox}[colback=white, colframe=black, sharp corners, boxrule=1pt, width=6.0cm, title=\textbf{Nutrition Facts (total)}, coltitle=black, colbacktitle=white, attach title to upper=\vspace{0.4em}]
        \small
        \hrule width \linewidth height 2pt \vspace{0.3em}
        \begin{tabular}{@{} l r l @{}}
        \noindent Energy        & \printValue{\totalKcal}    & kcal \\
        \noindent Fat           & \printValue{\totalFat}     & g    \\
        \noindent Saturated Fat & \printValue{\totalSatFat}  & g    \\
        \noindent Carbohydrate  & \printValue{\totalCarbs}   & g    \\
        \noindent Fiber         & \printValue{\totalFiber}   & g    \\
        \noindent Sugars        & \printValue{\totalSugar}   & g    \\
        \noindent Protein       & \printValue{\totalProtein} & g    \\
        \noindent Sodium / Salt & \printValue{\totalSalt}    & g    \\
        \end{tabular}
    \end{tcolorbox}%
}

% Updated box style that dynamically adapts its width to its surrounding container
\newtcolorbox{ingredientsboxfull}[1][\linewidth]{
    sharp corners, colback=gray!3, colframe=gray!30, boxrule=0.5pt,
    width=#1, top=10pt, bottom=10pt, left=10pt, right=10pt
}

% =======================================================
% CENTRAL DATA STORAGE: KEY-VALUE FORMAT (Values per 100g)
% =======================================================
\defIngredient{ryeflour}{name={Rye Whole Grain Flour}, kcal=339, protein=9, carbs=62.4, fat=2.3, satfat=0.3, sugar=5.64, fiber=14.8, salt=0.00}
\defIngredient{water}{name={Water}, kcal=0, protein=0, carbs=0, fat=0, satfat=0, sugar=0, fiber=0, salt=0}
\defIngredient{waterBoiling}{name={Water (boiling)}, kcal=0, protein=0, carbs=0, fat=0, satfat=0, sugar=0, fiber=0, salt=0}
\defIngredient{salt}{name={Sea Salt}, kcal=0, protein=0, carbs=0, fat=0, satfat=0, sugar=0, fiber=0, salt=99.0}
\defIngredient{sourdough}{name={Sourdough Starter}, kcal=169.5, protein=4.5, carbs=31.2, fat=1.15, satfat=0.15, sugar=2.82, fiber=7.4, salt=0.00}
\defIngredient{coarseRyeMeal75}{name={Coarse rye meal (grind size 7.5)}, kcal=339, protein=9, carbs=62.4, fat=2.3, satfat=0.3, sugar=5.64, fiber=14.8, salt=0.00}
\defIngredient{coarseRyeMeal75FootNote}{name={Coarse rye meal (grind size 7.5)\footnote{Setting on "hawos M001 M1 Grinder, 1.3 Kg, 360 W, Wood"}}, kcal=339, protein=9, carbs=62.4, fat=2.3, satfat=0.3, sugar=5.64, fiber=14.8, salt=0.00}
\defIngredient{coarseRyeMeal3}{name={Coarse rye meal (grind size 3)}, kcal=339, protein=9, carbs=62.4, fat=2.3, satfat=0.3, sugar=5.64, fiber=14.8, salt=0.00}


\begin{document}

\title{Meine Rezepte}
\author{}
\date{\today}
\maketitle

\tableofcontents
\clearpage

\section{Roggenschrotbrot}
\noindent\textsf{\small 1 Laib \textperiodcentered\ Dauer: ca. 3,5 Std. \hfill Brot \& Gebäck}
\vspace{1em}

\begin{ingredientsboxfull}
    \noindent
    % --- Linke Spalte: Brühstück ---
    \begin{minipage}[t]{0.45\linewidth}
        % r = amount aligned right, l = unit aligned left, l = ingridient aligned left
        \begin{tabular}{@{} r l l @{}}
            \multicolumn{3}{l}{\textbf{Brühstück:}} \\
            \midrule
            240,65 & g  & Schrot (Stufe 7.5) \\
            721,95 & ml & Wasser (kochend) \\
            48,13  & g  & Salz \\
            & &
        \end{tabular}
    \end{minipage}%
    \hfill
    \begin{minipage}[t]{0.45\linewidth}
        \begin{tabular}{@{} r l l @{}}
            \multicolumn{3}{l}{\textbf{Hauptteig:}} \\
            \midrule
            1347,63 & g  & Schrot (Stufe 7.5) \\
            828,23  & g  & Schrot (Stufe 3) \\
            1754,73 & ml & Wasser \\
            84,23   & g  & Sauerteigstarter \\
        \end{tabular}
    \end{minipage}
\end{ingredientsboxfull}

\vspace{1.5em}
\noindent\textbf{Zubereitung:}
\vspace{0.5em}

\begin{enumerate}[label=\textbf{\arabic*.}, leftmargin=1.8em, itemsep=1em, topsep=0pt]
    \item \textbf{Brühstück:} Zutaten für das Brühstück mit kochendem Wasser übergießen und klumpenfrei verrühren. Folie direkt auf die Oberfläche der Masse legen und für 8--10 Stunden verquellen lassen.
    
    \item \textbf{Sauerteigstarter Füttern 1:} Sauerteigstarter aus dem Kühlschrank nehmen und 1:1:1 füttern und (3--4 Stunden) warten, bis er sich verdoppelt hat. Sauerteigstarter nochmal 1:1:1 füttern und (3--4 Stunden) warten, bis er sich verdoppelt hat.
    
    \item \textbf{Sauerteigstarter Füttern 2:} Sauerteigstarter nochmal 1:1:1 füttern und (3--4 Stunden) warten, bis er sich verdoppelt hat.

     \item \textbf{Formen auskleiden:} Backpapier zurechtschneiden und Formen damit auslegen.
   
    \item \textbf{Hauptteig:} Sauerteigstarter in einer großen Schüssel mit der Wassermenge für den Hauptteig vermischen. Die restlichen Zutaten dazugeben und alles vermischen, bis alles homogen ist, und auf die Backformen aufteilen.
    
    \item \textbf{Gare:} Die Formen abdecken und 24 Stunden bei 18--20\,°C gehen lassen.
    
    \item \textbf{Backen:} wird auf der folgenden Seite fortgesetzt
    \begin{enumerate}
    	\item Backblech auf unterste Schiene in den Ofen einschieben und auf 250\,°C vorheizen.
    	\item Kastenformen mit Teig auf Gitter direkt über dem Backblech einschieben, 100\,ml kochendes Wasser auf das Blech kippen und den Ofen sofort schließen.
    	\item Nach 15 Minuten den Ofen zum Dampfablassen öffnen und das Backblech herausnehmen. Temperatur auf 200\,°C stellen.
    	\item Nach weiteren 65 Minuten ein Stechthermometer in die Brote stecken und weiterbacken, bis die Kerntemperatur mindestens 98\,°C erreicht.
    	\item Brote aus den Formen nehmen und 10 Minuten weiterbacken.
    	\item Ofen ausmachen, die Klappe leicht öffnen und die Brote mindestens 30 Minuten im Ofen abkühlen lassen.
    \end{enumerate}
         
\end{enumerate}

\vspace{2em}

\begin{tcolorbox}[colback=yellow!5, colframe=yellow!40, boxrule=0.5pt, sharp corners]
    \noindent\textbf{Interaktive time plan:} \\
    When isn the bread supposed to be done? \\
    Date (Format \texttt{DD.MM.YYYY}): 
    \Form{\TextField[name=endDate, width=6em, value={24.10.2026}, align=1]{}} 
    \hspace{1em}
    Time (Format \texttt{HH:MM}): 
    \Form{\TextField[name=endTime, width=4em, value={18:00}, align=1]{}}
\end{tcolorbox}

\vspace{1em}
\noindent\textbf{Calculated schedule:}
\vspace{0.5em}

\noindent
\begin{tabular}{r @{\hspace{1em}} c @{\hspace{1em}} p{\dimexpr\linewidth-10.5em\relax}}
    \timelineField{step1}{3610} & \textcolor{gray}{\textbullet} & \textbf{Brühstück:} Brühstück ansetzen und mit Folie abdecken\\
    \vphantom{M} & \textcolor{gray}{|} & \\
    \timelineField{step2}{3490} & \textcolor{gray}{\textbullet} & \textbf{Sauerteig auffrischen:} Sauerteigstarter füttern \\
    \vphantom{M} & \textcolor{gray}{|} & \\
    \timelineField{step3}{3250} & \textcolor{gray}{\textbullet} & \textbf{Sauerteig auffrischen:} Sauerteigstarter nochmal füttern\\
    \vphantom{M} & \textcolor{gray}{|} & \\
    \timelineField{step4}{3030}  & \textcolor{gray}{\textbullet} & \textbf{Backpapier:} Formen mit Backpapier auskleiden \\
    \vphantom{M} & \textcolor{gray}{|} & \\
    \timelineField{step5}{3010}  & \textcolor{gray}{\textbullet} & \textbf{Hauptteig:} Alles zusammen mischen und in Formen füllen \\
    \vphantom{M} & \textcolor{gray}{|} & \\
    \timelineField{step6}{3000}  & \textcolor{gray}{\textbullet} & \textbf{Lange Gare:} Formen abdecken und 24 Std. reifen lassen \\
    \vphantom{M} & \textcolor{gray}{|} & \\
    \timelineField{step7}{1585}   & \textcolor{gray}{\textbullet} & \textbf{Backofen vorheizen:} Ofen mit Blech auf 250\,°C bringen \\
    \vphantom{M} & \textcolor{gray}{|} & \\
    \timelineField{step8}{1560}   & \textcolor{gray}{\textbullet} & \textbf{Einschießen:} Brote mit Dampf backen (Temperatur 200\,°C) \\
    \vphantom{M} & \textcolor{gray}{|} & \\
    \timelineField{step9}{1470}   & \textcolor{gray}{\textbullet} & \textbf{Ofen abschalten:} Ofen ausmachen und leicht "offnen und Brote weitere 30 Minuten im Ofen lassen\\
    \vphantom{M} & \textcolor{gray}{|} & \\
    \timelineField{step10}{1440}    & \textcolor{gray}{\textbullet} & \textbf{Abk"uhlen:} Brote aus dem Ofen nehmen und abkühlen lassen. \\
    \vphantom{M} & \textcolor{gray}{|} & \\
    \timelineField{step11}{0}    & \textcolor{gray}{\textbullet} & \textbf{Fertig!:} Fertig zum Essen. \\
\end{tabular}


\end{document}
